\section{Robot Packs, Validation, Persistence, and Studio}
\label{sec:robotpacks}

\subsection{Robot Pack as the portable unit}

A \robotpack{} is a directory whose root document is
\pathcode{robot_pack.yaml}. It contains split semantic specifications, backend mapping
documents, and declared local assets. The current schema version is 0.1. The pack is the
unit a collaborator should copy, validate, open in Studio, and instantiate in a
runtime. It is not merely a pointer to a machine-local URDF tree.

\begin{lstlisting}[language={},caption={Representative Robot Pack layout.}]
smores_ep/
  robot_pack.yaml
  specs/
    module_types.yaml
    connector_types.yaml
    capabilities.yaml
  mappings/
    urdf_mapping.yaml
  assets/
    urdf/smores_ep.urdf
    mujoco/smores_ep.xml
    meshes/*.stl
    collision_meshes/*.obj
\end{lstlisting}

The root manifest defines identity, semantic version, description, bounded JSON-safe
metadata, model-view recipes, asset catalogs, split-document paths, and mappings. Split
documents keep frequently edited catalogs readable and give validation issues a
specific document and JSON-pointer-like path.

\subsection{Strict schema and stable identifiers}

Pydantic models are strict, frozen, and reject unknown fields. Coercion, non-finite
numbers, and invalid paths fail rather than being normalized silently. Semantic IDs use
lowercase snake case, begin with a letter, and are at most 64 characters. This is why a
connector type ID such as \code{EP} is invalid while \code{ep\_face} is valid. Display
names remain free-form and human-readable.

The distinction matters operationally. IDs appear in cross-references, runtime
instance IDs, mappings, graph keys, and events, so casual renaming is disruptive.
Names are presentation. Studio turns raw schema messages into this distinction and
suggests a normalized ID where possible.

\subsection{Module, joint, and connector semantics}

A module type selects a mechanical asset, identifies its root link, optionally records
mass, declares semantic joints and connector instances, and advertises capabilities.
It does not replicate all URDF links. A joint records source name, joint type,
parent/child link, normalized axis, exposed control modes, and unit-bearing limits.
The schema can express position, velocity, and effort; the current \mujoco{} adapter
implements effort only.

A connector instance includes:

\begin{itemize}
  \item stable ID and connector type reference;
  \item parent link;
  \item a named frame and/or numeric local pose;
  \item docking and approach axes expressed in the parent-link frame; and
  \item bounded custom metadata.
\end{itemize}

A connector type contains the shared mating contract: active flag, gender, mutual
compatibility catalog, discrete or continuous roll orientations, capture shape and
tolerances, physical connection intent, compliance/load data, undocking support,
docking policy, and metadata. Current connection intent vocabulary includes fixed,
compliant, hinge, ball, and custom; runtime support is narrower. This is deliberate:
the schema can preserve future intent while validation and capability checks refuse
unsupported execution.

\subsection{Custom metadata}

Metadata is intentionally JSON-compatible rather than an unbounded YAML escape hatch.
It is supported at the manifest, connector type, connector instance, and model-view
configuration levels. This allows platform provenance, calibration notes, UI hints,
and external identifiers without weakening standard fields. Metadata does not alter
docking semantics unless a registered consumer interprets it. The schema presently has
no arbitrary module-type metadata field, which prevents the report from overstating
the scope of the feature.

\subsection{Backend mappings and capabilities}

Mappings associate a semantic module with a selected asset and provide link, joint,
connector-frame, and constraint name translations. The design anticipates different
naming schemes in URDF, MJCF, and a future USD backend. Current \mujoco{} support does
not yet consume the full mapping vocabulary and generally expects semantic names to
match the selected asset. Named-frame-only connectors are refused; numeric local poses
are required for the current path.

Capabilities are either primitive actions or higher-level behaviors. They can require
connector types and can be advertised by modules. In version 0.1 they are validated
descriptions, not executable command signatures. This distinction prevents a pack from
claiming that a runtime behavior engine already exists.

\subsection{Safety and loading}

The loader uses safe YAML, rejects duplicate keys, checks schema version, and injects
catalog-key IDs before model validation. Pack and split-document symlinks are rejected.
Paths must be relative, normalized, contained under the pack root, and free of home,
absolute, Windows-drive, backslash, control-character, and parent-traversal forms.
Network asset retrieval is not permitted during load.

These restrictions serve two purposes. They make a pack genuinely portable, and they
prevent a downloaded data package from reading arbitrary files. Robot Packs also name
builders and backends by stable IDs; they do not contain executable Python import
paths.

\subsection{Validation profiles}

Validation produces deterministic, structured issues with severity, stable code,
document, path, entity, explanation, and suggested fix. The \emph{authoring} profile
permits incomplete data with warnings so a draft can be opened and refined. The
\emph{simulation} profile promotes missing connector limits, joint limits, or mappings
to errors where structural readiness requires them.

Simulation validation must not be confused with a solver run. It does not parse every
referenced URDF name, load a backend, prove collision quality, or guarantee controller
stability. It states that the semantic package is complete enough to attempt a
simulation.

\subsection{Canonical writing and atomic update}

The writer revalidates the in-memory model, renders deterministic split YAML, copies
declared assets without following symlinks or special entries, stages the result,
reloads it, and compares semantics before publication. A new \code{write} refuses an
existing destination. An \code{update} builds a sibling staging directory, renames the
original to a backup, publishes the new tree, and attempts rollback if publication
fails.

Canonical output intentionally discards comments, anchors, hand formatting, and unknown
files. This is the cost of treating the schema model, rather than arbitrary YAML text,
as authoring state. Studio's Save and Export commands use the same path, so CLI and GUI
do not generate subtly different packs.

\subsection{URDF import}

The dependency-light URDF importer caps input size, forbids symlinks, DTDs, and entity
declarations, and parses robot/link/joint trees, origins, axes, limits, masses, basic
geometry, mesh references, and material colors. It resolves relative, local
\code{file://}, and \code{package://} meshes against explicit roots; it never downloads
HTTP assets. The draft builder copies resolved meshes, rewrites URDF references,
creates an initial mapping, and derives module mass.

Connectors cannot be inferred reliably from ordinary URDF. Their semantics and precise
mating frames are authored afterward. Xacro expansion, transmissions, simulator
extensions, comprehensive OBJ/DAE sidecars, and texture copying remain future work.

\subsection{Robot Pack Studio}

Studio is a native PySide6 application with a PyVista/PyVistaQt viewport. Its Qt-free
\code{StudioProject} treats each edit functionally: dump the existing pack, apply a
bounded change, revalidate, and return a dirty replacement. This centralizes
referential integrity outside individual widgets.

The current interface supports:

\begin{itemize}
  \item opening a pack or creating one from URDF;
  \item project and URDF tree navigation;
  \item pack, connector-type, connector-instance, joint, and model-view property
        editing;
  \item connector-type add/edit/reference-safe removal;
  \item connector-instance add/edit/remove and parent-link association;
  \item row-oriented JSON metadata editing;
  \item authoring and simulation-profile validation;
  \item canonical YAML preview;
  \item atomic Save and non-overwriting Export As; and
  \item a per-launch local session log truncated at startup.
\end{itemize}

The viewport renders the URDF at zero joint configuration with detailed visual meshes,
material or fallback colors, smooth/PBR shading, multiple lights, shadows, an expanded
ground grid, link frames, joint axes, connector glyphs, and optional red wireframe
collision proxies. Selecting a link can initiate connector creation, and link selection
is highlighted.

\subsection{Studio limitations}

Studio currently visualizes one module at zero configuration. It does not animate
joints, render multi-module assemblies, pick connector glyphs in 3-D, provide a
connector transform gizmo, draw capture volumes, render named-frame-only connectors,
or show live contacts and forces. Texture images and normal maps are not bundled or
rendered. Rebuilding after an edit may reset the camera, and UI-level regression
coverage is lighter than core model coverage. These limitations motivate future
integration between the authoring and runtime shells but do not change the current
authority boundary.

